% 03-simulations-apprentissage.tex — État de l'art 2/5 : simulations
% interactives et apprentissage (≈ 7 pages).

\chapter{Simulations interactives et apprentissage}
\label{chap:simulations}

\section{Pédagogies actives en physique}

\begin{arediger}{1,5 page}
Engagement interactif contre enseignement traditionnel : gains normalisés
\parencite{hake1998}, méta-analyse de l'apprentissage actif \parencite{freeman2014}.
Ce qu'une simulation peut apporter dans un amphithéâtre ou en travail
autonome.
\end{arediger}

\section{Les simulations de physique : PhET et au-delà}

\begin{arediger}{2 pages}
PhET \parencite{wieman2008phet, adams2008phet} : principes de conception
(rétroaction immédiate, représentations multiples, exploration guidée),
études d'engagement. Comparaison laboratoire réel / virtuel
\parencite{finkelstein2005}. Tableau comparatif d'outils existants (PhET,
myPhysicsLab, GeoGebra, Algodoo…) : contrôle du temps, affichage des
forces, formules substituées, adaptation au cours. Montrer ce qui manque
pour un cours précis de BAC1.
\end{arediger}

\section{Apprentissage par découverte et guidage}

\begin{arediger}{1,5 page}
\textcite{dejong1998} : la découverte avec simulation n'est efficace qu'avec
du guidage (hypothèses, variables, consignes). Conséquence pour nos pages :
hypothèses affichées, verdict explicite (\enquote{le wagon fait le tour} /
\enquote{décollage à \(\theta_d\)}), défis du mode jeu.
\end{arediger}

\section{Charge cognitive et représentations multiples}

\begin{arediger}{2 pages}
Théorie de la charge cognitive \parencite{sweller1988} et principes du
multimédia \parencite{mayer2009} : contiguïté spatiale et temporelle,
signalement, cohérence. Justifie nos choix : formules à côté de la scène,
valeurs numériques substituées, code couleur constant des vecteurs, tout sur
un seul écran (section~\ref{sec:exercice-wagon}).
\end{arediger}
